\chapter{روش پیشنهادی}
\label{chp:proposed_method}
\section{مقدمه}
در این فصل به تشریح بهبودهای پیشنهادی بر روی دو الگوریتم افزایش وضوح تصویر شامل \lr{ScSR} و \lr{ZSSR} می‌پردازیم. هدف اصلی این بهبودها، کاهش چشمگیر زمان اجرا با حفظ حداکثری کیفیت بازسازی تصویر است. در بخش \ref{sec:scsr_improvement} بهبود پیشنهادی بر روی الگوریتم \lr{ScSR} ارائه می‌شود که شامل جایگزینی الگوریتم بهینه‌سازی در فاز بازنمایی تنک و بهینه‌سازی‌های پیاده‌سازی است. سپس در بخش \ref{sec:zssr_improvement} بهبودهای پیشنهادی بر روی الگوریتم \lr{ZSSR} تشریح می‌شود که شامل افزودن سیاست توقف زودهنگام و تنظیم ابرپارامترهای آموزش است. در هر بخش ابتدا مروری بر الگوریتم اصلی ارائه شده و سپس تغییرات پیشنهادی و انگیزه آن‌ها به تفصیل بیان می‌گردد.

\section{بهبود الگوریتم \lr{ScSR}}
\label{sec:scsr_improvement}

\subsection{مروری بر الگوریتم اصلی \lr{ScSR}}
الگوریتم \lr{ScSR}\LTRfootnote{Sparse Coding Super-Resolution} یک روش مبتنی بر بازنمایی تنک\LTRfootnote{Sparse Representation} برای افزایش وضوح تصویر است. این الگوریتم شامل دو فاز اصلی آموزش و اجرا می‌باشد.

\subsubsection{فاز آموزش}
در فاز آموزش، هدف یادگیری یک جفت دیکشنری وابسته برای فضاهای با وضوح بالا (\lr{HR}) و با وضوح پایین (\lr{LR}) است. مراحل این فاز به شرح زیر است:

\begin{enumerate}
\item \textbf{ساخت تصاویر با وضوح پایین:}
از تصاویر \lr{HR} با اعمال عملیات \lr{blur} و \lr{downsample} تصاویر \lr{LR} ساخته می‌شود. سپس تصاویر حاصل با روش درونیابی دوکعبی\LTRfootnote{Bicubic Interpolation} به اندازه اصلی بازگردانده می‌شوند.

\item \textbf{استخراج ویژگی:}
چهار فیلتر مشتقی بر روی تصاویر با وضوح پایین اعمال می‌شود تا چهار نقشه ویژگی\LTRfootnote{Feature Map} حاصل گردد.

\item \textbf{ساخت پچ‌ها:}
در هر موقعیت مکانی، از تصاویر \lr{HR} یک پچ استخراج شده و به بردار تبدیل می‌شود. همچنین از چهار نقشه ویژگی در همان موقعیت مکانی، چهار پچ استخراج و پس از الحاق\LTRfootnote{Concatenation} به یک بردار ویژگی واحد تبدیل می‌شوند.

\item \textbf{یادگیری دیکشنری:}
دو دیکشنری $D_h$ و $D_l$ به‌صورت همزمان با حل مسئله بهینه‌سازی زیر آموزش داده می‌شوند:
\begin{align}
D_h &= \arg\min_{D_h, Z} \|X_h - D_h Z\|_2^2 + \lambda \|Z\|_1 \\
D_l &= \arg\min_{D_l, Z} \|Y_l - D_l Z\|_2^2 + \lambda \|Z\|_1
\end{align}

برای سادگی محاسباتی، این دو مسئله با وزن‌دهی بر اساس تعداد ردیف‌ها به یک مسئله واحد تبدیل می‌شوند. فرآیند حل به‌صورت تکراری انجام می‌شود:
\begin{itemize}
\item ابتدا دیکشنری ثابت نگه داشته شده و ضرایب تنک با حل مسئله \lr{LASSO} کمینه می‌شوند.
\item سپس ضرایب ثابت نگه داشته شده و دیکشنری با حل مسئله \lr{QCQP}\LTRfootnote{Quadratically Constrained Quadratic Program} بهینه می‌شود.
\item این دو مرحله تا رسیدن به همگرایی تکرار می‌شوند.
\end{itemize}
\end{enumerate}

\subsubsection{فاز اجرا}
در فاز اجرا، برای افزایش وضوح یک تصویر ورودی مراحل زیر انجام می‌شود:

\begin{enumerate}
\item درونیابی دوکعبی تصویر ورودی تا رسیدن به اندازه \lr{HR}.
\item اعمال چهار فیلتر مشتقی و تولید چهار نقشه ویژگی.
\item پچ‌برداری به‌صورت \lr{raster-scan} با هم‌پوشانی\LTRfootnote{Overlap}.
\item الحاق چهار پچ از چهار نقشه ویژگی و تشکیل بردار ویژگی.
\item محاسبه میانگین هر پچ و کسر آن از مقادیر پیکسل‌ها به‌منظور بازسازی بر اساس بافت\LTRfootnote{Texture} و نه شدت\LTRfootnote{Intensity}.
\item \textbf{بازنمایی تنک محلی:} حل مسئله بهینه‌سازی برای یافتن ضرایب تنک $\alpha$:
\begin{align}
\alpha^* = \arg\min_{\alpha} \|F D_l \alpha - F y\|_2^2 + \lambda \|P D_h \alpha - w\|_2^2
\end{align}
که می‌توان آن را به فرم فشرده $\arg\min_{\alpha} \|D\alpha - y\|_2^2$ نوشت.
\item \textbf{قید بازسازی سراسری:} پس از بازسازی محلی تنک، قید بازسازی سراسری\LTRfootnote{Global Reconstruction Constraint} با حل مسئله $Y - SHX$ با استفاده از گرادیان کاهشی\LTRfootnote{Gradient Descent} اعمال می‌شود.
\end{enumerate}

\subsection{بهبودهای پیشنهادی بر الگوریتم \lr{ScSR}}
\label{subsec:scsr_proposed}

مهم‌ترین گلوگاه محاسباتی الگوریتم \lr{ScSR} در فاز اجرا، مرحله بازنمایی تنک محلی است. در الگوریتم اصلی، برای حل مسئله \lr{LASSO} از الگوریتم \lr{Feature-Sign Search} استفاده می‌شود که تا رسیدن به همگرایی ادامه می‌یابد و پیچیدگی زمانی آن $O(K^3)$ است، که $K$ تعداد اتم‌های دیکشنری را نشان می‌دهد. با توجه به اینکه مسئله بازنمایی تنک باید برای هر پچ تصویر به‌صورت مجزا حل شود، این پیچیدگی بالا منجر به زمان اجرای طولانی می‌گردد.

\subsubsection{جایگزینی با الگوریتم \lr{OMP}}

بهبود اصلی پیشنهادی، جایگزینی الگوریتم \lr{Feature-Sign Search} با الگوریتم \lr{Orthogonal Matching Pursuit (OMP)} است. برخلاف \lr{Feature-Sign Search} که تا رسیدن به همگرایی ادامه می‌یابد، الگوریتم \lr{OMP} تا تعداد مشخصی اتم (مثلاً ۴ یا ۶) پیش می‌رود و به‌صورت حریصانه\LTRfootnote{Greedy} در هر تکرار، اتم مناسب‌ را برای بازنمایی باقیمانده\LTRfootnote{Residual} انتخاب می‌کند.

مراحل الگوریتم \lr{OMP} به شرح زیر است. با فرض دیکشنری $D \in \mathbb{R}^{m \times K}$ و بردار مشاهده $y \in \mathbb{R}^m$:

\begin{enumerate}
\item \textbf{مقداردهی اولیه:}
باقیمانده اولیه $r_0 = y$ و مجموعه اتم‌های انتخابی $\Omega_0 = \emptyset$.

\item \textbf{محاسبه شباهت:}
در هر تکرار $k$، شباهت هر اتم $d_j$ با باقیمانده فعلی $r_{k-1}$ از طریق ضرب داخلی محاسبه می‌شود:
\begin{align}
h_j = |\langle d_j, r_{k-1} \rangle|, \quad j = 1, \ldots, K
\end{align}

\item \textbf{انتخاب اتم:}
اتمی که بیشترین شباهت را با باقیمانده دارد انتخاب و به مجموعه اتم‌ها اضافه می‌شود:
\begin{align}
j^* &= \arg\max_j h_j \\
\Omega_k &= \Omega_{k-1} \cup \{j^*\}
\end{align}

\item \textbf{محاسبه ضرایب:}
با استفاده از دیکشنری فرعی $D_{\Omega_k}$ (شامل اتم‌های انتخاب‌شده)، ضرایب $\alpha$ از طریق حل مسئله کمترین مربعات\LTRfootnote{Least Squares} محاسبه می‌شوند:
\begin{align}
\alpha_k = \arg\min_{\alpha} \|y - D_{\Omega_k} \alpha\|_2^2 = (D_{\Omega_k}^T D_{\Omega_k})^{-1} D_{\Omega_k}^T y
\end{align}

\item \textbf{به‌روزرسانی باقیمانده:}
\begin{align}
r_k = y - D_{\Omega_k} \alpha_k
\end{align}

\item \textbf{شرط توقف:}
الگوریتم تا رسیدن به تعداد اتم مدنظر (مثلاً $k_{\max} = 4$) یا کوچک شدن نرم باقیمانده از یک آستانه مشخص تکرار می‌شود.
\end{enumerate}

پیچیدگی زمانی الگوریتم \lr{OMP} برابر با $O(mkK)$ است که در آن $m$ اندازه بردار ویژگی (در اینجا $m = 36$)، $k$ تعداد اتم‌های انتخابی (در اینجا $k = 4$)، و $K$ تعداد کل اتم‌های دیکشنری است. با مقایسه $O(mkK)$ با $O(K^3)$ و با توجه به مقادیر $m = 36$ و $k = 4$، کاهش قابل توجه پیچیدگی محاسباتی آشکار می‌شود.

\subsubsection{بهینه‌سازی‌های پیاده‌سازی}

علاوه بر تغییر الگوریتم بازنمایی تنک، دو بهینه‌سازی پیاده‌سازی نیز اعمال شده است:

\begin{itemize}
\item \textbf{استخراج ویژگی یکباره:} در پیاده‌سازی اصلی، فیلترهای مشتقی به‌صورت مجزا برای هر پچ اعمال می‌شدند. در روش پیشنهادی، مطابق با متن مقاله اصلی، استخراج ویژگی یک‌بار برای کل تصویر انجام شده و سپس پچ‌ها از نقشه‌های ویژگی محاسبه‌شده استخراج می‌شوند. این تغییر از محاسبات تکراری جلوگیری می‌کند.

\item \textbf{استفاده از \lr{float32} به‌جای \lr{float64}:} نوع داده اعشاری از دقت مضاعف ۶۴ بیتی به دقت تکی ۳۲ بیتی تغییر یافته است. این تغییر ضمن حفظ دقت کافی برای کاربرد افزایش وضوح تصویر، مصرف حافظه را نصف کرده و سرعت محاسبات عددی را افزایش می‌دهد.
\end{itemize}


\section{بهبود الگوریتم \lr{ZSSR}}
\label{sec:zssr_improvement}

\subsection{مروری بر الگوریتم اصلی \lr{ZSSR}}
الگوریتم \lr{ZSSR}\LTRfootnote{Zero-Shot Super-Resolution} یک روش افزایش وضوح تصویر بدون ناظر\LTRfootnote{Zero-Shot} است که برای هر تصویر تست، یک شبکه عصبی مجزا آموزش می‌دهد. اصول کلیدی این الگوریتم عبارتند از:

\subsubsection{شبکه مخصوص تصویر}
برای هر تصویر تست $I$، یک شبکه عصبی تمام‌ کانولوشنی\LTRfootnote{Fully Convolutional} به‌صورت خودنظارتی\LTRfootnote{Self-Supervised} آموزش داده می‌شود. از خود تصویر $I$ با عمل \lr{downsampling} با ضریب مقیاس $s$، نسخه کوچک‌تر $I{\downarrow s}$ ساخته می‌شود. شبکه آموزش می‌بیند تا $I$ را از $I{\downarrow s}$ بازسازی کند و پس از آموزش، همان شبکه بر روی $I$ اعمال می‌شود تا خروجی $I{\uparrow s}$ تولید گردد.

\subsubsection{ساخت داده آموزشی از یک تصویر}
از آنجا که تنها یک تصویر در دسترس است، مجموعه داده آموزشی به‌صورت داخلی و چندمقیاسی ساخته می‌شود. چند نسخه کوچک‌تر از تصویر اصلی ($I_0, I_1, I_2, \ldots, I_n$) به‌عنوان \lr{HR fathers} تولید شده و هر یک با ضریب مقیاس ثابت $s$ کاهش وضوح داده می‌شوند تا \lr{LR sons} متناظر ساخته شوند. برای افزایش تعداد نمونه‌ها، از تقویت داده\LTRfootnote{Data Augmentation} هندسی شامل ۴ چرخش و بازتاب‌ها\LTRfootnote{Flip} استفاده می‌شود که داده‌ها را تا ۸ برابر افزایش می‌دهد.

\subsubsection{افزایش وضوح تدریجی}
برای ضرایب مقیاس بزرگ، فرآیند افزایش وضوح به‌صورت چندمرحله‌ای انجام می‌شود. در هر مرحله خروجی \lr{SR} تولید شده و نسخه‌های کاهش‌یافته آن به مجموعه \lr{HR fathers} اضافه می‌گردد. این چرخه تا رسیدن به ضریب مقیاس نهایی ادامه می‌یابد.

\subsubsection{معماری شبکه و تنظیمات آموزش}
شبکه شامل ۸ لایه پنهان با ۶۴ کانال و تابع فعال‌سازی \lr{ReLU} است. ورودی ابتدا با درونیابی به اندازه خروجی بزرگ شده و شبکه تنها باقیمانده\LTRfootnote{Residual} بین تصویر درونیابی‌شده و تصویر \lr{HR} را یاد می‌گیرد. تابع هزینه \lr{L1}، بهینه‌ساز \lr{Adam} و نرخ یادگیری با شروع از $10^{-3}$ و کاهش تدریجی تا $10^{-6}$ استفاده می‌شود. در هر تکرار یک برش تصادفی\LTRfootnote{Random Crop} با اندازه ثابت $128 \times 128$ از یک جفت \lr{father–son} انتخاب شده و \lr{fathers} بزرگ‌تر با احتمال بیشتری انتخاب می‌شوند.

\subsubsection{پس‌تصویرسازی و خودترکیب}
الگوریتم از تکنیک خودترکیب هندسی\LTRfootnote{Self-Ensemble} استفاده می‌کند: ۸ حالت چرخش و بازتاب بر تصویر اعمال شده، برای هر تبدیل افزایش وضوح انجام می‌گیرد، پس‌تصویرسازی\LTRfootnote{Back-Projection} اعمال شده، همه هم‌ترازسازی\LTRfootnote{Align} می‌شوند و میانه\LTRfootnote{Median} گرفته می‌شود. سپس یک‌بار دیگر پس‌تصویرسازی بر تصویر نهایی حاصل از میانه اعمال می‌گردد.

\subsubsection{سیاست نرخ یادگیری}
در الگوریتم اصلی، نرخ یادگیری بر اساس شیب خطای بازسازی در بازه مشخصی از تکرارهای گذشته بررسی می‌شود. هر \lr{run\_test\_every} تکرار، تصویر \lr{SR} بر مبنای بزرگ‌ترین \lr{LR Son} ساخته شده و با \lr{HR Father} مقایسه می‌شود تا میانگین مربعات خطا (\lr{MSE}) محاسبه گردد. این \lr{MSE} هم برای ارزیابی پیشرفت یادگیری و هم برای تصمیم‌گیری درباره کاهش نرخ یادگیری استفاده می‌شود. هر \lr{learning\_rate\_policy\_check\_every} تکرار، شیب \lr{MSE} در \lr{learning\_rate\_slope\_range} تکرار گذشته بررسی شده و در صورت عدم بهبود کافی، نرخ یادگیری کاهش می‌یابد. الگوریتم زمانی متوقف می‌شود که نرخ یادگیری به $9 \times 10^{-6}$ برسد.

\subsection{بهبودهای پیشنهادی بر الگوریتم \lr{ZSSR}}
\label{subsec:zssr_proposed}

با توجه به اینکه الگوریتم \lr{ZSSR} برای هر تصویر تست یک شبکه مجزا آموزش می‌دهد، زمان اجرای آن به‌صورت مستقیم وابسته به تعداد تکرارهای آموزش است. هدف بهبودهای پیشنهادی، کاهش تعداد تکرارهای غیرضروری و در نتیجه کاهش زمان اجرا با حداقل افت کیفیت است.

\subsubsection{سیاست توقف زودهنگام مبتنی بر \lr{MSE} بازسازی}
\label{subsec:early_stop}

مهم‌ترین بهبود پیشنهادی، افزودن یک سیاست توقف زودهنگام\LTRfootnote{Early Stopping} مبتنی بر \lr{MSE} بازسازی است. در الگوریتم اصلی، \lr{MSE} بازسازی تنها برای تنظیم نرخ یادگیری استفاده می‌شد. در روش پیشنهادی، این معیار به‌عنوان شاخصی برای تشخیص اشباع یادگیری\LTRfootnote{Learning Saturation} نیز به‌کار گرفته می‌شود.

سازوکار توقف زودهنگام به این صورت عمل می‌کند: در هر آزمون (که هر \lr{run\_test\_every} تکرار انجام می‌شود)، \lr{MSE} فعلی ($\text{MSE}_{\text{current}}$) با بهترین \lr{MSE} مشاهده‌شده ($\text{MSE}_{\text{best}}$) مقایسه می‌شود. اگر بهبود معنادار، یعنی اختلاف بیشتر از $10^{-6}$، مشاهده نشود، شمارنده عدم بهبود افزایش می‌یابد:

\begin{align}
\text{no\_improvement\_count} = 
\begin{cases}
\text{no\_improvement\_count} + 1 & \text{اگر } \text{MSE}_{\text{best}} - \text{MSE}_{\text{current}} < 10^{-6} \\
0 & \text{در غیر این صورت}
\end{cases}
\end{align}

در صورتی که $\text{MSE}_{\text{current}} < \text{MSE}_{\text{best}}$، مقدار $\text{MSE}_{\text{best}}$ به‌روزرسانی می‌شود. اگر در ۳ آزمون متوالی بهبود معنادار مشاهده نشود ($\text{no\_improvement\_count} \geq 3$)، آموزش متوقف می‌گردد.

انتخاب آستانه $10^{-6}$ به این دلیل صورت گرفته که تغییرات \lr{MSE} کوچک‌تر از این مقدار تأثیر قابل ادراکی بر کیفیت بصری تصویر نهایی ندارند و ادامه آموزش پس از این نقطه عملاً بازدهی محاسباتی نخواهد داشت. انتخاب ۳ آزمون متوالی نیز برای اطمینان از اینکه عدم بهبود موقتی نیست و واقعاً اشباع یادگیری رخ داده است، صورت گرفته است.

\subsubsection{تنظیم ابرپارامترهای آموزش}

سه تغییر دیگر در ابرپارامترهای آموزش اعمال شده است:

\begin{enumerate}
\item \textbf{کاهش اندازه برش تصادفی:}
اندازه برش تصادفی (\lr{crop\_size}) از $128 \times 128$ به $64 \times 64$ کاهش یافته است. این تغییر حجم محاسبات هر تکرار را کاهش می‌دهد. با توجه به اینکه الگوریتم \lr{ZSSR} از تکرارشوندگی الگوها در مقیاس‌های مختلف تصویر بهره می‌برد، پچ‌های کوچک‌تر نیز می‌توانند الگوهای تکراری داخلی تصویر را به‌خوبی در بر گیرند.

\item \textbf{کاهش فاصله آزمون و بررسی نرخ یادگیری:}
پارامترهای \lr{run\_test\_every} و \lr{learning\_rate\_policy\_check\_every} از ۵۰ به ۴۰ تکرار کاهش یافته‌اند. این تغییر باعث می‌شود ارزیابی پیشرفت یادگیری و تصمیم‌گیری درباره کاهش نرخ یادگیری با فواصل کوتاه‌تری انجام شود. در نتیجه، اشباع یادگیری زودتر تشخیص داده شده و سیاست توقف زودهنگام سریع‌تر فعال می‌گردد.

\item \textbf{کاهش بازه بررسی شیب نرخ یادگیری:}
پارامتر \lr{learning\_rate\_slope\_range} از ۲۵۰ به ۲۰۰ تکرار کاهش یافته است. این پارامتر تعداد تکرارهای گذشته‌ای را تعیین می‌کند که برای محاسبه شیب \lr{MSE} و تصمیم‌گیری درباره کاهش نرخ یادگیری مورد استفاده قرار می‌گیرند. همچنین حداقل تعداد تکرار لازم برای شروع بررسی نرخ یادگیری نیز به همین مقدار کاهش یافته است. کاهش این بازه به واکنش سریع‌تر سیاست نرخ یادگیری به تغییرات روند یادگیری منجر می‌شود.
\end{enumerate}

\subsection{تأثیر ترکیبی بهبودها}

بهبودهای پیشنهادی بر الگوریتم \lr{ZSSR} به‌صورت مکمل عمل می‌کنند. سیاست توقف زودهنگام سازوکار اصلی کاهش تکرارهای غیرضروری را فراهم می‌آورد. کاهش فاصله آزمون و بازه بررسی نرخ یادگیری، حساسیت تشخیص اشباع یادگیری را افزایش می‌دهد و کاهش اندازه برش تصادفی، هزینه محاسباتی هر تکرار را کم می‌کند. ترکیب این تغییرات در مجموع منجر به کاهش چشمگیر زمان اجرا می‌شود. جزئیات نتایج حاصل از این بهبودها در فصل \ref{chp:results} ارائه خواهد شد.
